\chapter{El espacio de funciones integrables}\label{cap:v2-08}

La integral de $|f-g|$ mide el error total de una aproximación. Para convertirla en distancia hay que identificar funciones iguales casi en todas partes. Esa identificación permite estudiar simultáneamente completitud, aproximación y paso al límite. Reunimos las lecturas 13, 16 y 17 de Valencia \cite{valencia2}; el argumento por series para completitud también se explica en las notas de Maximenko \cite{maximenkoLp}.

\section{Seminormas, cocientes e igualdad casi en todas partes}\label{lectura:v2-13}
\begin{definicion}
Una \emph{seminorma}\index{seminorma} en un espacio vectorial $V$ sobre $\F=\R$ o $\C$ es una aplicación $p:V\to[0,\infty)$ tal que $p(\lambda v)=|\lambda|p(v)$ y $p(u+v)\le p(u)+p(v)$. Es norma si además $p(v)=0$ implica $v=0$.
\end{definicion}
Las condiciones implican $p(0)=0$ y $|p(u)-p(v)|\le p(u-v)$. El conjunto $W=\{v:p(v)=0\}$ es un subespacio: la homogeneidad y la desigualdad triangular preservan el valor cero. En el cociente $V/W$, dos vectores representan la misma clase si su diferencia pertenece a $W$. Las operaciones $[u]+[v]=[u+v]$, $\lambda[v]=[\lambda v]$ no dependen de los representantes, porque $W$ es subespacio. La fórmula
\[\|[v]\|=p(v)
\]
está bien definida por la desigualdad inversa; es una norma porque una clase tiene norma cero exactamente cuando es la clase de $W$.

\begin{ejemplo}
En las funciones Riemann integrables sobre $[a,b]$, $p(f)=\int_a^b|f|$ es una seminorma: cambiar el valor en un punto no cambia la integral. Si $X$ es compacto y $A\subset X$, la restricción $C(X)\to C(A)$ tiene núcleo $W_A=\{f:f|_A=0\}$, de modo que $C(X)/W_A$ identifica funciones con la misma restricción. Si $X$ no es compacto, la expresión $\sup_X|f|$ puede ser infinita para $f\in C(X)$; para usarla como norma hay que trabajar en $C_b(X)$ o imponer compacidad. Las afirmaciones sobre ideales de $C(X)$ tampoco se trasladan sin hipótesis del caso compacto al general.
\end{ejemplo}

Sea $\mathscr L^1(\mu)$ el espacio de funciones medibles $f$ con $\int|f|<\infty$. Las propiedades de la integral muestran que es un espacio vectorial y que $p(f)=\int|f|$ es seminorma. Su núcleo está formado exactamente por las funciones nulas casi en todas partes. Definimos
\[L^1(\mu)=\mathscr L^1(\mu)/\{f:f=0\text{ casi en todas partes}\},\qquad
\|f\|_1=\int_X|f|\,d\mu.
\]
Por comodidad escribimos $f$ para su clase. Una afirmación puntual sobre elementos de $L^1$ exige escoger representantes; una afirmación de norma debe ser independiente de esa elección.

\begin{ejercicio}[Lectura 13, ejercicio 1]\label{ej:v2-l13-1}
Si $T:V\to U$ es lineal, pruebe $V/\ker T\cong\operatorname{im}T$.
\end{ejercicio}

\begin{ejercicio}[Lectura 13, ejercicio 2]\label{ej:v2-l13-2}
Demuestre $\dim V=\dim W+\dim(V/W)$ para un subespacio $W$.
\end{ejercicio}

\begin{ejercicio}[Lectura 13, ejercicio 3]\label{ej:v2-l13-3}
Si $W\subset\ker T$, construya la aplicación inducida $V/W\to U$ y pruebe su unicidad.
\end{ejercicio}


\section{Completitud y aproximación por funciones simples}\label{lectura:v2-16}
\begin{teorema}[Riesz--Fischer para $L^1$]\label{teo:v2-L1-Banach}\index{Riesz--Fischer, teorema de}\index{L1@$L^1$!completitud}
El espacio $L^1(\mu)$ es de Banach para cualquier espacio de medida.
\end{teorema}
\begin{proof}
Primero consideremos una serie $\sum_ku_k$ tal que $\sum_k\|u_k\|_1<\infty$. Elegimos representantes medibles y ponemos $H=\sum_k|u_k|$. La convergencia monótona da
\[\int H\,d\mu=\sum_k\int|u_k|\,d\mu<\infty.
\]
En consecuencia, $H$ es finita fuera de un conjunto nulo. Allí la serie converge absolutamente a una función $u$, que definimos como cero en la excepción. Es medible, $|u|\le H$ casi en todas partes y
\[\left\|u-\sum_{k=1}^nu_k\right\|_1\le\sum_{k>n}\|u_k\|_1\longrightarrow0.
\]
Así, toda serie absolutamente convergente en norma converge en $L^1$.

Ahora sea $(f_n)$ de Cauchy. Escogemos una subsucesión con $\|f_{n_{k+1}}-f_{n_k}\|_1<2^{-k}$. La serie de diferencias converge en $L^1$ por el argumento anterior, y sus sumas parciales telescópicas muestran que $f_{n_k}$ converge en $L^1$ a cierto $f$. Dado $\eps>0$, tomamos un índice de Cauchy $N$ con $\|f_n-f_m\|_1<\eps/2$ para $m,n\ge N$, y un $n_k\ge N$ con $\|f_{n_k}-f\|_1<\eps/2$. Para $n\ge N$, la desigualdad triangular da $\|f_n-f\|_1<\eps$. La sucesión completa converge.
\end{proof}

\begin{proposicion}\label{prop:v2-simples-densas}\index{densidad!de funciones simples}
Las funciones simples integrables son densas en $L^1(\mu)$.
\end{proposicion}
\begin{proof}
Si $f\ge0$ es integrable, escojamos simples $0\le s_n\uparrow f$. Entonces $\|f-s_n\|_1=\int f-\int s_n\to0$ por convergencia monótona; cada $s_n$ es integrable por comparación. Para $f$ real aplicamos el argumento a sus partes positiva y negativa. Para $f$ compleja lo aplicamos a las partes real e imaginaria y usamos $|u+iv|\le|u|+|v|$. En una representación canónica de una simple integrable, todo conjunto con coeficiente no nulo tiene medida finita, pues $|a|\mu(E)\le\int|s|<\infty$.
\end{proof}

\section{Separabilidad del espacio de medida}
Para $A,B\in\Acal$, pongamos $\rho(A,B)=\mu(A\mathbin\triangle B)$, permitiendo infinito, y
\[d_\mu(A,B)=\frac{\rho(A,B)}{1+\rho(A,B)},\qquad \frac{\infty}{1+\infty}:=1.
\]
Esta es una seudométrica. En efecto, $A\triangle C\subset(A\triangle B)\cup(B\triangle C)$, así que $\rho$ satisface la desigualdad triangular extendida. La función creciente $h(t)=t/(1+t)$ es subaditiva: para $s,t<\infty$,
\[h(s+t)=\frac{s+t}{1+s+t}\le\frac{s}{1+s}+\frac{t}{1+t}.
\]
Si un término es infinito, la desigualdad es inmediata porque el miembro derecho es al menos uno. Identificando conjuntos que difieren en un nulo se obtiene un espacio métrico, llamado aquí \emph{álgebra de medida}. Decimos que el espacio de medida es separable si esta métrica admite un conjunto numerable denso. La condición $d_\mu(A,B)<\eps/(1+\eps)$ equivale a $\mu(A\triangle B)<\eps$; no deben confundirse las dos cantidades.

\begin{teorema}\label{teo:v2-medida-separable}\index{separabilidad!del algebra de medida@del álgebra de medida}
Si el álgebra de medida es separable, entonces $L^1(\mu)$ es separable. Si además $\mu$ es semifinita, entonces es $\sigma$-finita.
\end{teorema}
\begin{proof}
Todo subespacio de un espacio métrico separable es separable: las bolas con centros en un conjunto numerable denso y radios racionales forman una base numerable; de cada elemento de la base que interseca el subespacio elegimos un punto del subespacio. Apliquemos esto al conjunto de clases representadas por conjuntos de medida finita. Obtenemos una familia numerable $(E_j)$, con $\mu(E_j)<\infty$, que aproxima en medida de diferencia simétrica a cualquier conjunto de medida finita.

Las combinaciones lineales finitas de los $\ind{E_j}$ con coeficientes racionales, o racionales complejos, forman un conjunto numerable. Aproximan a cada simple integrable: si $s=\sum_{i=1}^ra_i\ind{A_i}$ con $\mu(A_i)<\infty$, elegimos $E_{j_i}$ con $\sum_i|a_i|\mu(A_i\triangle E_{j_i})$ tan pequeño como se desee. Luego aproximamos cada $a_i$ por un racional $q_i$, haciendo pequeña $\sum_i|a_i-q_i|\mu(E_{j_i})$. Ambas estimaciones se suman por la desigualdad triangular. La proposición \ref{prop:v2-simples-densas} concluye separabilidad de $L^1$.

Para la última afirmación, recordemos que semifinitud significa que todo conjunto de medida infinita contiene un subconjunto medible de medida \emph{positiva y finita}; equivalentemente, la medida de cualquier conjunto es el supremo de las medidas finitas de sus subconjuntos. Sea $E=\bigcup_jE_j$. Si $X\setminus E$ contuviera un conjunto $B$ de medida positiva finita, entonces $\mu(B\triangle E_j)\ge\mu(B)>0$ para todo $j$, contradiciendo la densidad de $(E_j)$ entre los conjuntos de medida finita. Por semifinitud, $\mu(X\setminus E)=0$. Así, los $E_j$ y el complemento nulo constituyen una cubierta numerable por conjuntos de medida finita.
\end{proof}

\begin{observacion}
La separabilidad de $L^1$ no implica separabilidad del álgebra de medida, como muestra el ejercicio \ref{ej:v2-l16-4}. Tampoco la medida finita implica separabilidad de $L^1$. Para verlo, sea $I$ no numerable y $X=\{0,1\}^I$. Consideremos los conjuntos $\pi_J^{-1}(B)$, donde $J\subset I$ es numerable y $B$ es boreliano en $\{0,1\}^J$. Forman una $\sigma$-álgebra: una sucesión de tales conjuntos depende sólo de la unión numerable de sus coordenadas. Sobre cada espacio numerable de coordenadas existe la probabilidad de monedas independientes, construida con los dígitos binarios de un punto uniforme de $[0,1)$; un cilindro que prescribe $r$ dígitos tiene medida $2^{-r}$, y la unicidad de extensión determina la medida de todos los borelianos. Definimos $\mu(\pi_J^{-1}(B))$ mediante esa probabilidad. La definición es independiente de $J$, porque al pasar a la unión de dos conjuntos de coordenadas las marginales coinciden en cilindros y, por unicidad, en todos los borelianos. La aditividad numerable se comprueba llevando toda la sucesión a un único conjunto numerable de coordenadas. Tenemos así una probabilidad en $X$. Para $i\ne j$, las funciones coordenadas $x\mapsto x_i$ y $x\mapsto x_j$ distan $1/2$ en $L^1$. Una familia no numerable con separación uniforme positiva impide separabilidad: bolas disjuntas de radio menor que $1/4$ requerirían puntos distintos de cualquier conjunto denso.
\end{observacion}

\section{Ejercicios sobre cocientes, densidad y separabilidad}
\begin{ejercicio}[Lectura 16, ejercicio 1]\label{ej:v2-l16-1}
Identifique la completación de las funciones Riemann integrables en $[a,b]$ con la seminorma integral.
\end{ejercicio}

\begin{ejercicio}[Lectura 16, ejercicio 2]\label{ej:v2-l16-2}
Si $d$ es una seudométrica, pruebe que identificar puntos a distancia cero produce una métrica.
\end{ejercicio}

\begin{ejercicio}[Lectura 16, ejercicio 3]\label{ej:v2-l16-3}
Caracterice la separabilidad del cociente mediante una sucesión de representantes.
\end{ejercicio}

\begin{ejercicio}[Lectura 16, ejercicio 4]\label{ej:v2-l16-4}
Compare la separabilidad del álgebra de medida de $\N$ con medida de conteo y la de su espacio $L^1$.
\end{ejercicio}

\begin{ejercicio}[Lectura 16, ejercicio 5]\label{ej:v2-l16-5}
Pruebe que $L^1(\R,m)$ es separable.
\end{ejercicio}

\begin{ejercicio}[Lectura 16, ejercicio 6]\label{ej:v2-l16-6}
Examine la afirmación de que la medida de conteo en la $\sigma$-álgebra numerable--conumerable de $\R$ produce $L^1$ separable sin ser $\sigma$-finita. Dé una versión correcta.
\end{ejercicio}


\section{Convergencia en media y control de la masa}\label{lectura:v2-17}
\begin{definicion}
Decimos que $f_n\to f$ \emph{en media}\index{convergencia!en media} si $\|f_n-f\|_1\to0$.
\end{definicion}
Por Markov,
\[\mu\{|f_n-f|>a\}\le\frac{\|f_n-f\|_1}{a},\qquad a>0.
\]
Así, convergencia en media implica convergencia en medida y permite extraer una subsucesión que converge casi uniformemente y casi en todas partes. La desigualdad tiene dirección $\le$, que es esencial para esta deducción. Si $\mu(X)<\infty$, la convergencia uniforme de funciones integrables implica convergencia en media: primero $|f|\le|f_N|+1$ para algún $N$ muestra que $f$ es integrable, y luego $\|f_n-f\|_1\le\mu(X)\sup_X|f_n-f|$.

\begin{teorema}[Convergencia dominada en medida]\label{teo:v2-dominada-medida}
Si $f_n\to f$ en medida y $|f_n|\le g\in L^1$, entonces $f\in L^1$ y $\|f_n-f\|_1\to0$.
\end{teorema}
\begin{proof}
Una subsucesión converge a $f$ casi en todas partes, por lo que $|f|\le g$ casi en todas partes y $f$ es integrable. Si la conclusión fallara, existirían $\eps_0>0$ y una subsucesión con $\|f_{n_k}-f\|_1\ge\eps_0$ para todo $k$. Esa subsucesión aún converge en medida y admite otra que converge casi en todas partes. En esta última, $|f_{n_{k_j}}-f|\le2g$; convergencia dominada da normas tendiendo a cero, contradicción.
\end{proof}

\begin{lema}\label{lema:v2-control-integral}
Para $g\in L^1$, la integral de $|g|$ es absolutamente continua respecto de la medida: dado $\eps>0$, existe $\delta>0$ tal que $\mu(E)<\delta$ implica $\int_E|g|<\eps$. Además, existe un conjunto $A$ de medida finita con $\int_{X\setminus A}|g|<\eps$.
\end{lema}
\begin{proof}
Las funciones $|g|\ind{\{|g|>M\}}$ tienden a cero casi en todas partes cuando $M$ crece. Convergencia dominada permite elegir $M>0$ para que su integral sea menor que $\eps/2$. Tomemos $\delta=\eps/(2M)$. Si $\mu(E)<\delta$, entonces
\[\int_E|g|\le M\mu(E)+\int_{\{|g|>M\}}|g|<\eps.\]
Para la segunda afirmación tomamos $A_k=\{|g|\ge1/k\}$. Markov da $\mu(A_k)\le k\|g\|_1<\infty$, mientras $|g|\ind{X\setminus A_k}\to0$ puntualmente y está dominada por $|g|$. La convergencia dominada permite elegir $k$.
\end{proof}

\begin{teorema}[Criterio de Vitali en $L^1$]\label{teo:v2-vitali-L1}\index{Vitali!criterio de convergencia}
Sean $f_n,f\in L^1(\mu)$. Entonces $\|f_n-f\|_1\to0$ si y sólo si se cumplen:
\begin{enumerate}
\item $f_n\to f$ en medida;
\item para cada $\eps>0$ existe $A\in\Acal$ con $\mu(A)<\infty$ y $\sup_n\int_{X\setminus A}|f_n|<\eps$;
\item para cada $\eps>0$ existe $\delta>0$ tal que $\mu(E)<\delta$ implica $\sup_n\int_E|f_n|<\eps$.
\end{enumerate}
\end{teorema}
\begin{proof}
Si hay convergencia en $L^1$, ya probamos (i). Para (ii), elija $N$ con $\|f_n-f\|_1<\eps/2$ para $n\ge N$. Por el lema anterior hay un conjunto de medida finita que recoge toda la integral de $|f|$ salvo $\eps/2$. Añadiendo un conjunto de medida finita apropiado para cada una de las funciones $f_1,\ldots,f_{N-1}$, obtenemos una unión finita $A$ que sirve a todas. Para $n\ge N$, la integral exterior es menor que $\|f_n-f\|_1+\int_{X\setminus A}|f|<\eps$; para los primeros índices se impuso directamente. Para (iii), escogemos $N$ con la misma estimación y aplicamos la primera parte del lema a la familia finita $f,f_1,\ldots,f_{N-1}$; el mínimo de sus tolerancias sirve a toda la sucesión, usando de nuevo la desigualdad triangular.

Recíprocamente, fijemos $\eps>0$. Por (ii), y ampliando $A$ mediante el lema aplicado a $f$, podemos lograr $\mu(A)<\infty$ y
\[\sup_n\int_{X\setminus A}|f_n|<\eps,\qquad \int_{X\setminus A}|f|<\eps.
\]
Por (iii) y la continuidad absoluta de la integral de $|f|$, escogemos $\delta>0$ tal que cualquier $E$ de medida menor que $\delta$ satisfaga $\sup_n\int_E|f_n|<\eps$ y $\int_E|f|<\eps$. Sea $a=\eps/(1+\mu(A))>0$. Por (i), desde cierto índice $E_n=\{|f_n-f|>a\}$ tiene medida menor que $\delta$. Descomponiendo $X$ en $X\setminus A$, $A\cap E_n$ y $A\setminus E_n$, obtenemos
\[\|f_n-f\|_1<2\eps+2\eps+a\mu(A)<5\eps.
\]
Como $\eps$ es arbitrario, la norma tiende a cero.
\end{proof}
\begin{observacion}[Una hipótesis que no puede omitirse]
La condición $\mu(A)<\infty$ en (ii) es necesaria. Sin ella se podría tomar $A=X$, haciendo vacía la condición. En $\R$, $f_n=n^{-1}\ind{[0,n]}$ converge uniformemente a cero y cumple (iii), pues $\int_E|f_n|\le m(E)$. Sin embargo, $\|f_n\|_1=1$. La masa se dispersa sobre conjuntos cada vez mayores. En contraste, $n\ind{(0,1/n)}$ concentra su masa en conjuntos pequeños y viola (iii). Los dos controles de Vitali impiden fenómenos diferentes.
\end{observacion}

\section{Ejercicios sobre integrabilidad, convexidad y traslaciones}
\begin{ejercicio}[Lectura 17, ejercicio 1]\label{ej:v2-l17-1}
Demuestre la unicidad del límite en media.
\end{ejercicio}

\begin{ejercicio}[Lectura 17, ejercicio 2]\label{ej:v2-l17-2}
Analice $f_n=n\ind{[1/n,2/n]}$ en $[0,1]$.
\end{ejercicio}

\begin{ejercicio}[Lectura 17, ejercicio 3]\label{ej:v2-l17-3}
Si $f_n\to f$ en medida y $|f_n|\le g\in L^1$, pruebe $\int f_n\to\int f$.
\end{ejercicio}

\begin{ejercicio}[Lectura 17, ejercicio 4]\label{ej:v2-l17-4}
Compare las condiciones de control de masa de Vitali cuando $|f_n|\le g\in L^1$.
\end{ejercicio}

\begin{ejercicio}[Lectura 17, ejercicio 5: Jensen]\label{ej:v2-l17-5}\index{Jensen, desigualdad de}
Sea $\mu(X)=1$, $f\in L^1$ real y $\varphi:\R\to\R$ convexa. Pruebe continuidad de $\varphi$, existencia de una recta de soporte en cada punto y $\varphi(\int f)\le\int\varphi(f)$.
\end{ejercicio}

\begin{ejercicio}[Lectura 17, ejercicio 6]\label{ej:v2-l17-6}
Si $\mu(X)<\infty$, $\varphi$ es continua y $|\varphi(t)|\le K|t|$ para $|t|\ge K$, demuestre que $f\in L^1$ real implica $\varphi\circ f\in L^1$. Establezca la versión para medida infinita.
\end{ejercicio}

\begin{ejercicio}[Lectura 17, ejercicio 7]\label{ej:v2-l17-7}
Bajo las hipótesis del ejercicio \ref{ej:v2-l17-6}, pruebe que $f_n\to f$ en $L^1$ implica $\varphi(f_n)\to\varphi(f)$ en $L^1$.
\end{ejercicio}

\begin{ejercicio}[Lectura 17, ejercicio 8]\label{ej:v2-l17-8}
Si $f\in L^1(\R)$ y $f_x(t)=f(t-x)$, pruebe que $x\mapsto f_x$ es continua con valores en $L^1$.
\end{ejercicio}
